% GENERATED FILE. Edit canonical lessons, then run npm run generate:books.
% canonical-source-hash: fnv1a64:e3f33031347fb5e1
% canonical-lessons: GE-C317-erzaehlung, GE-C317-nacherzaehlen, GE-C317-vorlesen, GE-C317-ueberliefern, GE-C317-schicksal, GE-R317-review-310-317

\chapter{Weitererzählen}
\label{ch:ge-weitererzahlen}
\hlchaptermodality{317}

\begin{chapteropening}
\textbf{By the end of this chapter:} \emph{I can talk about telling and passing on stories: a narrative, to retell, to read aloud, to hand down and fate.}

Complete the last lesson of chapter 317: I can talk about telling and passing on stories: a narrative, to retell, to read aloud, to hand down and fate.
\end{chapteropening}

\section[die Erzählung]{die Erzählung — the narrative, the short story}

\label{lesson:GE-C317-erzaehlung}



\begin{quote}
Before the new one: say the German for the chapter, then the German for the sequel.
\end{quote}

\subsection*{You'll want to know: die Erzählung}
\textbf{die Erzählung} — \textquotedblleft{}the narrative, the short story\textquotedblright{}.

\begin{quote}
Die Erzählung spielt in einem kleinen Dorf.
\end{quote}

\begin{quote}
\emph{The narrative is set in a small village.}
\end{quote}

Plural: \emph{die Erzählungen}. As a book, a short story.

\begin{cousinweb}
\emph{Erzählung} is built on \emph{erzählen} \textquotedblleft{}to tell\textquotedblright{}, from \emph{zählen} \textquotedblleft{}to count\textquotedblright{}, Old High German \emph{zellen} (EtymDB), with the ending \emph{-ung}.
\end{cousinweb}

\begin{grammarlens}[title={What you've built}]
One more word for telling, explaining and describing at B1.
\end{grammarlens}

\subsection*{Your turn}
\begin{itemize}
\raggedright
  \item \emph{Say it:} \emph{die Erzählung}
  \item \emph{Say it:} \emph{die Erzählung}, once more
  \item \emph{Say it:} a sentence of your own with \emph{die Erzählung}
  \item \emph{Recall:} say the German for the chapter, then the German for the sequel, then say \emph{die Erzählung} again
\end{itemize}

\begin{tcolorbox}[breakable,colback=teal!4,colframe=teal!35!black,title={Before you move on}]
What does die Erzählung mean? (\textquotedblleft{}The narrative, the short story\textquotedblright{}.) Say it once more.
\end{tcolorbox}
\section[nacherzählen]{nacherzählen — to retell}

\label{lesson:GE-C317-nacherzaehlen}



\begin{quote}
Before the new one: say the German for the sequel, then the German for the narrative.
\end{quote}

\subsection*{You'll want to know: nacherzählen}
\textbf{nacherzählen} — \textquotedblleft{}to retell\textquotedblright{}.

\begin{quote}
Die Schüler erzählten die Geschichte nach.
\end{quote}

\begin{quote}
\emph{The pupils retold the story.}
\end{quote}

Separable: \emph{sie erzählt nach}, \emph{sie hat nacherzählt}.

\begin{cousinweb}
\emph{nacherzählen} joins \emph{nach} \textquotedblleft{}after\textquotedblright{} and \emph{erzählen} \textquotedblleft{}to tell\textquotedblright{}, built on \emph{zählen} \textquotedblleft{}to count\textquotedblright{}, Proto-Germanic \emph{taljaną} (EtymDB): to tell after someone.
\end{cousinweb}

\begin{grammarlens}[title={What you've built}]
One more word for telling, explaining and describing at B1.
\end{grammarlens}

\subsection*{Your turn}
\begin{itemize}
\raggedright
  \item \emph{Say it:} \emph{nacherzählen}
  \item \emph{Say it:} \emph{nacherzählen}, once more
  \item \emph{Say it:} a sentence of your own with \emph{nacherzählen}
  \item \emph{Recall:} say the German for the sequel, then the German for the narrative, then say \emph{nacherzählen} again
\end{itemize}

\begin{tcolorbox}[breakable,colback=teal!4,colframe=teal!35!black,title={Before you move on}]
What does nacherzählen mean? (\textquotedblleft{}To retell\textquotedblright{}.) Say it once more.
\end{tcolorbox}
\section[vorlesen]{vorlesen — to read aloud}

\label{lesson:GE-C317-vorlesen}



\begin{quote}
Before the new one: say the German for the narrative, then the German for to retell.
\end{quote}

\subsection*{You'll want to know: vorlesen}
\textbf{vorlesen} — \textquotedblleft{}to read aloud\textquotedblright{}.

\begin{quote}
Der Vater las den Kindern eine Geschichte vor.
\end{quote}

\begin{quote}
\emph{The father read the children a story.}
\end{quote}

Separable, strong: \emph{liest vor}, \emph{las vor}, \emph{hat vorgelesen}.

\begin{cousinweb}
\emph{vorlesen} joins \emph{vor} \textquotedblleft{}before, in front of\textquotedblright{} and \emph{lesen} \textquotedblleft{}to read\textquotedblright{}, from Old High German \emph{lesan}, Proto-Germanic \emph{lesaną} \textquotedblleft{}to gather\textquotedblright{} (EtymDB): first gathering, then reading.
\end{cousinweb}

\begin{grammarlens}[title={What you've built}]
One more word for telling, explaining and describing at B1.
\end{grammarlens}

\subsection*{Your turn}
\begin{itemize}
\raggedright
  \item \emph{Say it:} \emph{vorlesen}
  \item \emph{Say it:} \emph{vorlesen}, once more
  \item \emph{Say it:} a sentence of your own with \emph{vorlesen}
  \item \emph{Recall:} say the German for the narrative, then the German for to retell, then say \emph{vorlesen} again
\end{itemize}

\begin{tcolorbox}[breakable,colback=teal!4,colframe=teal!35!black,title={Before you move on}]
What does vorlesen mean? (\textquotedblleft{}To read aloud\textquotedblright{}.) Say it once more.
\end{tcolorbox}
\section[überliefern]{überliefern — to hand down, to pass on}

\label{lesson:GE-C317-ueberliefern}



\begin{quote}
Before the new one: say the German for to retell, then the German for to read aloud.
\end{quote}

\subsection*{You'll want to know: überliefern}
\textbf{überliefern} — \textquotedblleft{}to hand down, to pass on\textquotedblright{}.

\begin{quote}
Die Lieder wurden von den Großeltern überliefert.
\end{quote}

\begin{quote}
\emph{The songs were handed down by the grandparents.}
\end{quote}

Often passive: \emph{Es ist überliefert} — it has come down to us.

\begin{cousinweb}
\emph{überliefern} is built on \emph{liefern} \textquotedblleft{}to deliver\textquotedblright{}, Middle High German \emph{lieveren} (EtymDB), which is usually traced through French \emph{livrer} to Latin \emph{līberāre} \textquotedblleft{}to set free\textquotedblright{}.
\end{cousinweb}

\begin{grammarlens}[title={What you've built}]
One more word for telling, explaining and describing at B1.
\end{grammarlens}

\subsection*{Your turn}
\begin{itemize}
\raggedright
  \item \emph{Say it:} \emph{überliefern}
  \item \emph{Say it:} \emph{überliefern}, once more
  \item \emph{Say it:} a sentence of your own with \emph{überliefern}
  \item \emph{Recall:} say the German for to retell, then the German for to read aloud, then say \emph{überliefern} again
\end{itemize}

\begin{tcolorbox}[breakable,colback=teal!4,colframe=teal!35!black,title={Before you move on}]
What does überliefern mean? (\textquotedblleft{}To hand down, to pass on\textquotedblright{}.) Say it once more.
\end{tcolorbox}
\section[das Schicksal]{das Schicksal — the fate, destiny}

\label{lesson:GE-C317-schicksal}



\begin{quote}
Before the new one: say the German for to read aloud, then the German for to hand down.
\end{quote}

\subsection*{You'll want to know: das Schicksal}
\textbf{das Schicksal} — \textquotedblleft{}the fate, destiny\textquotedblright{}.

\begin{quote}
Sie konnte ihr Schicksal nicht ändern.
\end{quote}

\begin{quote}
\emph{She could not change her fate.}
\end{quote}

Plural: \emph{die Schicksale}.

\begin{cousinweb}
\emph{Schicksal} is built on \emph{schicken}, which in Middle High German meant \textquotedblleft{}to arrange, to order\textquotedblright{} (EtymDB), with the ending \emph{-sal}: what has been arranged for you.
\end{cousinweb}

\begin{grammarlens}[title={What you've built}]
One more word for telling, explaining and describing at B1.
\end{grammarlens}

\subsection*{Your turn}
\begin{itemize}
\raggedright
  \item \emph{Say it:} \emph{das Schicksal}
  \item \emph{Say it:} \emph{das Schicksal}, once more
  \item \emph{Say it:} a sentence of your own with \emph{das Schicksal}
  \item \emph{Recall:} say the German for to read aloud, then the German for to hand down, then say \emph{das Schicksal} again
\end{itemize}

\begin{tcolorbox}[breakable,colback=teal!4,colframe=teal!35!black,title={Before you move on}]
What does das Schicksal mean? (\textquotedblleft{}The fate, destiny\textquotedblright{}.) Say it once more.
\end{tcolorbox}
\section[(dialogue)]{Words from chapters 310–317 — a review}

\label{lesson:GE-R317-review-310-317}



\begin{quote}
Say the German for the carnival and the Shrovetide carnival.
\end{quote}

\subsection*{Your turn}
Say these aloud:

\begin{itemize}
\raggedright
  \item the carnival, the Shrovetide carnival, the costume, to dress up, the custom
  \item the Christmas market, the Christmas tree, the Advent wreath, the giving of Christmas presents, the gingerbread
  \item the public festival, the funfair, the merry-go-round, the traditional dress, the band
  \item the fairy tale, the legend, the fable, the witch, the wizard
  \item to conjure, to bewitch, to turn into, the curse, the dragon
  \item the knight, the king, the princess, the giant, the dwarf
  \item the narrator, the plot, the character, the chapter, the sequel
  \item the narrative, to retell, to read aloud, to hand down, the fate
\end{itemize}

\begin{tcolorbox}[breakable,colback=teal!4,colframe=teal!35!black,title={Before you move on}]
The carnival? (\textbf{der Karneval}.) The Shrovetide carnival? (\textbf{der Fasching}.) The costume? (\textbf{das Kostüm}.) To dress up? (\textbf{verkleiden}.) The custom? (\textbf{der Brauch}.) The Christmas market? (\textbf{der Weihnachtsmarkt}.) The Christmas tree? (\textbf{der Weihnachtsbaum}.) The Advent wreath? (\textbf{der Adventskranz}.) The giving of Christmas presents? (\textbf{die Bescherung}.) The gingerbread? (\textbf{der Lebkuchen}.) The public festival? (\textbf{das Volksfest}.) The funfair? (\textbf{der Jahrmarkt}.) The merry-go-round? (\textbf{das Karussell}.) The traditional dress? (\textbf{die Tracht}.) The band? (\textbf{die Kapelle}.) The fairy tale? (\textbf{das Märchen}.) The legend? (\textbf{die Legende}.) The fable? (\textbf{die Fabel}.) The witch? (\textbf{die Hexe}.) The wizard? (\textbf{der Zauberer}.) To conjure? (\textbf{zaubern}.) To bewitch? (\textbf{verzaubern}.) To turn into? (\textbf{verwandeln}.) The curse? (\textbf{der Fluch}.) The dragon? (\textbf{der Drache}.) The knight? (\textbf{der Ritter}.) The king? (\textbf{der König}.) The princess? (\textbf{die Prinzessin}.) The giant? (\textbf{der Riese}.) The dwarf? (\textbf{der Zwerg}.) The narrator? (\textbf{der Erzähler}.) The plot? (\textbf{die Handlung}.) The character? (\textbf{die Figur}.) The chapter? (\textbf{das Kapitel}.) The sequel? (\textbf{die Fortsetzung}.) The narrative? (\textbf{die Erzählung}.) To retell? (\textbf{nacherzählen}.) To read aloud? (\textbf{vorlesen}.) To hand down? (\textbf{überliefern}.) The fate? (\textbf{das Schicksal}.)
\end{tcolorbox}
